\subsection{标准群的滤过}

我们再次采用定义 1 的记号。选取一个数 \(a > 1\) 和 K 的实赋值 \(v\)\(\mathbf{K}\)，使得对所有 \(|x| = a^{-v(x)}\) 都有 \(x \in \mathbf{K}\)\(a\)（交换代数，第 VI 章，§6，命题 3）。若 a 是 A 的包含于 m 中的非零（因而是开的）理想，令 \(\mathrm{G}(a)\)\(a\) 表示坐标属于 a 的 G 中元素的集合。若 \(\lambda \in \mathbf{R}\)，令 \(a_{\lambda}\)（分别地，\(a_{\lambda}^{+}\)）表示满足 \(x \in \mathbf{K}\)（分别地，\(v(x) \geqslant \lambda\)）的 \(v(x) > \lambda\) 的集合；于是 \(a_0 = A\)，\(a_0^+ = m\)。对于 \(x = (x_1, \ldots, x_r) \in G\)，记

\[
\omega (x) = \inf (v (x _ {1}), \dots , v (x _ {r})).\tag{5}
\]

命题 5. 设 G 是标准群。

\begin{enumerate}
\def\labelenumi{(\roman{enumi})}
\item
  若 \(\mathfrak{a}\) 是包含于 \(\mathbf{A}\) 中的 \(\mathfrak{m}\) 的非零理想，则 \(\mathbf{G}(\mathfrak{a})\) 是 \(\mathbf{G}\) 的开正规子群。
\item
  当 \(\mathbf{G}(\mathfrak{a}_{\lambda})\) 时，\(\lambda > 0\)\(e\) 构成 \(\mathbf{G}\) 中 e 的邻域基。
\item
  假设 \(\mathfrak{a}_{\lambda} \subset \mathfrak{a}\)（当 \(\lambda \geqslant \lambda_0\) 时），并给 \(G(\mathfrak{a}) / G(\mathfrak{a}_{\lambda})\) 赋予离散拓扑（当 \(\lambda \geqslant \lambda_0\) 时）。则拓扑群 \(G(\mathfrak{a})\) 是各群 \(G(\mathfrak{a}) / G(\mathfrak{a}_{\lambda})\) 的逆极限。
\item
  设 \(\mathfrak{a},\mathfrak{b}\) 是包含于 \(\mathfrak{m}\) 中的 A 的非零理想，满足 \(\mathfrak{a} \supset \mathfrak{b} \supset \mathfrak{a}^2\)。从 G(a) 到 (a/b) × · · · × (a/b) 的映射 \(x \mapsto (x_1 \bmod \mathfrak{b}, \ldots, x_r \bmod \mathfrak{b})\) 在取商后定义出从 G(a)/G(b) 到加法群 (a/b) × · · · × (a/b) 的同构。
\end{enumerate}

若 \(x \in G\) 且 \(y \in G(a)\)\(x\)\(x.y\)\(a\)，则 x 和 x.y 的坐标模 a 相等。因此，对于 G 中的 \(x'\)、\(x''\)\(G\) 以及 G(a) 中的 \(y'\)、\(y''\)\(G(a)\)\(x'.x''\)\((x'.y')\)\((x''.y'')\)\(a\)，x'.x'' 与 (x'.y') . (x''.y'') 的坐标模 a 相等。这证明了 (i)。

\begin{enumerate}
\def\labelenumi{(\roman{enumi})}
\setcounter{enumi}{1}
\item
  显然。
\item
  由上可知，这由一般拓扑，第 III 章，§ 7，命题 2 得到。
\end{enumerate}

若 \(x \in G(a)\) 且 \(y \in G(a)\)\(x, y\)\(x + y\)，则由 § 5 的公式 (4)，x、y 的坐标与 x + y 的坐标模 \(G(a^2)\) 同余。这证明了 (iv)。

推论。假设 K 局部紧，令 \(q = \operatorname{Card}(\mathbf{A} / \mathfrak{m})\)。

\begin{enumerate}
\def\labelenumi{(\roman{enumi})}
\item
  若 \(a = m^{a}\) 且 \(b = m^{b}\)，其中 \(b \geqslant a \geqslant 1\)，则 \(\mathrm{G}(a)/\mathrm{G}(b)\) 是基数为 \(q^{r(b-a)}\) 的 p-群。
\item
  G(a) 是 p-群的逆极限。
\end{enumerate}

\(\mathrm{G}(\mathfrak{a})/\mathrm{G}(\mathfrak{b})\) 中元素的数目是 \((\mathrm{Card}(\mathfrak{a}/\mathfrak{b}))^{r}\)；若 \(b=a+1\)，则 a/b 是 A/m 上的一维向量空间，从而在此情形下得到 (i)；一般情形由对 b-a 作归纳得到。断言 (ii) 由 (i) 和命题 5 (iii) 得到。

命题 6. 设 \(\mathfrak{a}, \mathfrak{b}, \mathfrak{c}, \mathfrak{c}'\) 是包含于 \(\mathbf{A}\) 中的 \(\mathfrak{m}\) 的非零理想，满足

\[
\mathfrak {c} ^ {\prime} \subset \mathfrak {c}, \quad \mathfrak {a b} \subset \mathfrak {c}, \quad \mathfrak {a b} ^ {2} \subset \mathfrak {c} ^ {\prime}, \quad \mathfrak {a} ^ {2} \mathfrak {b} \subset \mathfrak {c} ^ {\prime}.
\]

若 \(x \in \mathrm{G}(\mathfrak{a})\) 且 \(y \in \mathrm{G}(\mathfrak{b})\)，则 \(x^{[-1]}y^{[-1]}x.y, x.y.x^{[-1]}y^{[-1]}\)、 和 \([x,y]\) 属于 \(\mathrm{G}(\mathfrak{c})\)，并且模 \(\mathrm{G}(\mathfrak{c}')\) 同余。

由 § 5，第 2 号，命题 1，存在 \(c_{\alpha \beta} \in \mathbf{A}^r\)，使得

\[
x ^ {[ - 1 ]}. y ^ {[ - 1 ]}. x. y - [ x, y ] = \sum_ {| \alpha | + | \beta | \geqslant 3} c _ {\alpha \beta} x ^ {\alpha} y ^ {\beta}.
\]

若 \(x = 0\) 或 \(y = 0\)，则 \(x^{[-1]}y^{[-1]}x.y - [x,y] = 0\)；因此 \(c_{0\beta} = c_{\alpha 0} = 0\)。另一方面，条件

\[
x \in \mathrm{G} (\mathfrak {a}), \quad y \in \mathrm{G} (\mathfrak {b}), \quad | \alpha | \geqslant 1, \quad | \beta | \geqslant 1, \quad | \alpha | + | \beta | \geqslant 3
\]

蕴含

\[
c _ {\alpha \beta} x ^ {\alpha} y ^ {\beta} \in \mathrm{G} (a ^ {2} b + a b ^ {2}) \subset \mathrm{G} (c ^ {\prime})
\]

因而 \(x^{[-1]} \cdot y^{[-1]}x \cdot y - [x, y] \in \mathrm{G}(\mathfrak{c}')\)。类似地可得

\[
x. y. x ^ {[ - 1 ]}. y ^ {[ - 1 ]} - [ x, y ] \in \mathrm{G} (\mathfrak {c} ^ {\prime}).
\]

最后，由 § 5，公式 (13)，\([x,y] \in G(\mathfrak{ab}) \subset G(\mathfrak{c})\)。

命题 7. (i) 族 \((\mathrm{G}(\mathfrak{a}_{\lambda}))\) 是 \(\mathbf{G}\) 上的中心滤过（第二章，§ 4，第 4 号，定义 2）。

\[
\text { (ii) } For \lambda \in \mathbf {R} _ {+} ^ {*}, G (a _ {\lambda}) = \{x \in G | \omega (x) \geqslant \lambda \}, G (a _ {\lambda} ^ {+}) = \{x \in G | \omega (x) > \lambda \}.
\]

\begin{enumerate}
\def\labelenumi{(\roman{enumi})}
\setcounter{enumi}{1}
\tightlist
\item
  显然。我们证明 (i)。显然 \(G(a_{\lambda}) = \bigcap_{\mu < \lambda} G(a_{\mu})\)，且 \(G = \bigcup_{\lambda > 0} G(a_{\lambda})\)。另一方面，若 \(x \in G(a_{\lambda})\) 且 \(y \in G(a_{\mu})\)，则
\end{enumerate}

\[
x ^ {[ - 1 ]}. y ^ {[ - 1 ]}. x. y \in G (a _ {\lambda + \mu})
\]

将命题 6 应用于 \(a = a_{\lambda}, b = a_{\mu}, c = c' = a_{\lambda + \mu}\) 得到。

由第二章，§ 4，第 4 号，我们可以构造与带中心滤过 \((G(a_\lambda))\) 的群 G 相联系的群 gr(G)。对所有 \(G_\lambda = G(a_\lambda)/G(a_\lambda^+)\) 写 \(\lambda > 0\)，得到 \(gr(G) = \bigoplus_{\lambda>0} G_\lambda\)。回忆一下（同上，命题 1），G 中的换位子使我们能够在 gr(G) 上定义括号，使 gr(G) 成为李代数，具体如下：若 \(x \in G_\lambda\) 且 \(y \in G_\mu\)，取 \(\bar{x}\) 在 \(G(a_\lambda)\) 中的一个代表元 x，以及 \(\bar{y}\) 在 \(G(a_\mu)\) 中的一个代表元 y；则 \([\bar{x},\bar{y}]\) 是 \(x^{[-1]}.y^{[-1]}.x.y \in G(a_{\lambda+\mu})\) 在 \(G_{\lambda+\mu}\) 中的类。由命题 6（取 \(a = a_\lambda, b = a_\mu, c = a_{\lambda+\mu}, c' = a_{\lambda+\mu}\)），可见 \([\bar{x},\bar{y}]\) 也是 \([x,y]\) 在 \(G_{\lambda+\mu}\) 中的类。因此，当把 G 看作由 \(K^r\)\(G(a_\lambda)\) 滤过的李子代数  时，相关分次李代数（第二章，§ 4，第 3 号）等于 gr(G)。

